\section{Conclusion}
\label{sec:conclusion}

This article established a formally verified property calculus for
finite integer lists represented by recursive head--tail decomposition.

The core proved properties can be summarized as follows, for lists
$L,A,B,P,S \in \mathbb{L}$, values $x,e,v \in \mathbb{S}$, and valid
natural indices.

\begin{gather*}
\begin{aligned}
|L| > 0
&\implies L_{|L|-1} = \operatorname{last}(L)
&&\text{[Last Element Identity]} \\
0 < i < |L|
&\implies L_i = \operatorname{tail}(L)_{i-1}
&&\text{[Tail Access Shift]} \\
0 \leq f < t < |L|
&\implies L[f \dots t]
 = L[f \dots (t - 1)] \mathbin{\texttt{++}} (L_t :: L_e)
&&\text{[Slice Append Consistency]}
\end{aligned}
\end{gather*}

\begin{gather*}
\begin{aligned}
\operatorname{sum}(L)
&= \sum_{i=0}^{|L|-1} L_i
&&\text{[Sum Matches Summation]} \\
\operatorname{sum}(x :: L)
&= x + \operatorname{sum}(L)
&&\text{[Left Append Preserves Sum]} \\
\operatorname{sum}(A \mathbin{\texttt{++}} B)
&= \operatorname{sum}(A) + \operatorname{sum}(B)
&&\text{[Sum over Concatenation]} \\
\operatorname{sum}(A \mathbin{\texttt{++}} B)
&= \operatorname{sum}(B \mathbin{\texttt{++}} A)
&&\text{[Commutativity of Sum]} \\
(\forall x \in L,\ x > 0) \land L \neq L_e
&\implies \operatorname{sum}(L) > 0
&&\text{[Sum Positivity]}
\end{aligned}
\end{gather*}

\begin{gather*}
\begin{aligned}
\operatorname{product}(x :: L_e)
&= x
&&\text{[Singleton Product]} \\
\operatorname{product}(A \mathbin{\texttt{++}} (e :: B))
&= e \cdot \operatorname{product}(A \mathbin{\texttt{++}} B)
&&\text{[Product Pull-Out Element]} \\
\operatorname{product}(A \mathbin{\texttt{++}} B)
&= \operatorname{product}(A) \cdot \operatorname{product}(B)
&&\text{[Product over Concatenation]} \\
\operatorname{product}(A \mathbin{\texttt{++}} B)
&= \operatorname{product}(B \mathbin{\texttt{++}} A)
&&\text{[Commutativity of Product]} \\
(\forall x \in L,\ x > 0)
&\implies \operatorname{product}(L) > 0
&&\text{[Positive Product]}
\end{aligned}
\end{gather*}

\begin{gather*}
L \neq L_e \land (\forall x \in L,\ x > 0)
\implies \operatorname{product}(L) \bmod \operatorname{head}(L) = 0
\quad \text{[Head Divides Product]} \\
(\forall x \in L,\ x > 0)
\implies \forall x \in L,\ \operatorname{product}(L) \bmod x = 0
\quad \text{[All Elements Divide Product]} \\
\begin{aligned}
&e > 0 \land (\forall x \in P,\ x > 0) \land (\forall x \in S,\ x > 0) \\
&\implies \operatorname{product}(P \mathbin{\texttt{++}} (e :: S)) \bmod e = 0
\quad \text{[Inserted Element Divides Product]}
\end{aligned}
\end{gather*}

\begin{equation*}
\begin{aligned}
&(\forall x \in L,\ x > v) \land 0 \leq i < |L|
  \implies L_i > v
  &&\text{[Bound at Index]} \\
&(\forall x \in A,\ x > v) \land (\forall x \in B,\ x > v) \\
&\qquad \implies \forall x \in A \mathbin{\texttt{++}} B,\ x > v
  &&\text{[Bound over Concatenation]} \\
&(\forall x \in L,\ x > v) \land 0 \leq k \leq |L| \\
&\qquad \implies (\forall x \in \text{front},\ x > v)
     \land (\forall x \in \text{back},\ x > v)
  &&\text{[Split Preserves Lower Bound]} \\
&(\forall x \in A,\ x < b) \land (\forall x \in B,\ x < b) \\
&\qquad \implies \forall x \in A \mathbin{\texttt{++}} B,\ x < b
  &&\text{[Bound over Concatenation, Upper]} \\
&(\forall x \in L,\ x < b) \land 0 \leq k \leq |L| \\
&\qquad \implies (\forall x \in \text{front},\ x < b)
     \land (\forall x \in \text{back},\ x < b)
  &&\text{[Split Preserves Upper Bound]} \\
&\operatorname{slice}(L,f,t)
  = \operatorname{headRecursiveSlice}(L,f,t) \\
&\qquad = \operatorname{indexRangeValues}(L,f,t)
  &&\text{[Slice Equivalence]}
\end{aligned}
\end{equation*}

\begin{gather*}
\operatorname{value}_{h,G}(i + 1) - \operatorname{value}_{h,G}(i)
= G_i
\quad \text{[Adjacent Difference]} \\
\begin{aligned}
&\operatorname{value}_{\operatorname{shift}(h,G)}(i + 1)
 - \operatorname{value}_{\operatorname{shift}(h,G)}(i) \\
&= \operatorname{value}_{h,G}(i + 2) - \operatorname{value}_{h,G}(i + 1)
\quad \text{[Gap Translation]}
\end{aligned} \\
\operatorname{rotateAt}(L,k).\text{contains}(x)
\iff L.\text{contains}(x)
\quad \text{[Rotation Same Elements]} \\
|\operatorname{rotateAt}(L,k)| = |L|
\quad \text{[Rotation Same Size]} \\
\operatorname{sum}(\operatorname{rotateAt}(L,k)) = \operatorname{sum}(L)
\quad \text{[Rotation Same Sum]}
\end{gather*}

All of these properties are verified in the source references cited
throughout the article. Appendix~A collects the Scala excerpts that are
useful to keep near the text; each excerpt links back to its maintained
source file.

Collectively, these results give one verified calculus for decomposing
and composing finite lists, aggregating and bounding their values,
relating their elements to list products, preserving period while
tracking adjacent values and gaps under shifted-list transformations, and
preserving membership, size, sum, and element bounds under rotation.

\section{Future Work}
\label{sec:future-work}

Extending lists via integration (cumulative sums) and derivation (gap
extraction) would formalize two dual operations that map between a list
and its accumulated or decomposed form. These operations connect the
finite list algebra presented here to the theory of discrete sequences
and differences.

Two related list disciplines have verified source proofs but are not
developed here as headline properties: maintaining ascending order under
insertion and filtering (\texttt{SortedList}), and preserving a lower or
upper numeric bound under filtering alone rather than the append/split
operations covered in Section~8 (\texttt{MinBoundList},
\texttt{MaxBoundList}). A dedicated treatment of ordered and
bounded-filter list variants is left as future work.
